\documentclass[12pt, ukrainian]{article}
\usepackage[a4paper]{geometry}
\setlength\parindent{0mm}
\geometry{top=1.0cm,bottom=1.0cm,left=1.3cm,right=1.3cm}
\pagestyle{empty}
\usepackage[T2A]{fontenc}
\usepackage[utf8]{inputenc}
\usepackage[ukrainian]{babel}
\usepackage{graphicx}
\usepackage{caption}
\usepackage{caption}
\DeclareCaptionFormat{nocaption}{}
\captionsetup{format=nocaption,aboveskip=0pt,belowskip=0pt}
\usepackage[Export]{adjustbox}
\adjustboxset{max size={0.9\linewidth}{0.9\paperheight}}
\usepackage{verbatim, listings, clrscode3e, fancyvrb}
\def\code#1{\Verb!#1!}
\begin{document}
%begin
\begin {large}

\textbf{\_\_.11.2025 Контрольна робота, варіант  \No { 4 }}\par


{

\begin{center}
\textbf{Завдання 1}
\end{center}
Нехай int var=13; Один потік збільшив змінну var на 3, а інший збільшив її в три рази.

гонка даних (data race), стан гонки (race condition)

Виберіть усі правильні твердження (якщо такі є).\par
\vspace{2mm}
(\,1\,) Тут гарантовано нема гонки даних, але може бути стан гонки.%bad

(\,2\,) Тут може бути тільки стан гонки.%bad

(\,3\,) Тут гарантовано є стан гонки.%good

(\,4\,) Тут може не бути стану гонки.%bad



\begin{center}
\textbf{Завдання 2}
\end{center}
Виберіть усі правильні твердження (якщо такі є).\par
\vspace{2mm}

(\,1\,) Функція async запускає виконання в новому потоці.%bad

(\,2\,) У стані гонки програма відчайдушно намагається захопити ресурси процесора.%bad

(\,3\,) Секція є критичною, якщо під час її виконання не можна генерувати винятки.%bad

(\,4\,) Коли потік спить, він не споживає ресурс процесора.%good



\begin{center}
\textbf{Завдання 3}
\end{center}
Виберіть усі правильні твердження (якщо такі є).\par
\vspace{2mm}

(\,1\,) Вільних від очікувань структур даних не існує.%bad

(\,2\,) Операція над структурою даних без блокувань завжди виконається за обмежений час.%bad

(\,3\,) Якщо весь код структури даних утворює одну критичну секцію, то така структура даних може використовуватись в конкурентній програмі, але по-справжньо\-му конкурентною не є.%good

(\,4\,) На операціях структури даних без блокувань взаємоблокувань потоків виникнути не може.%bad



\begin{center}
\textbf{Завдання 4}
\end{center}
Виберіть усі правильні твердження (якщо такі є).\par
\vspace{2mm}

(\,1\,) Автоматичні змінні не можуть розділятися між потоками.%bad

(\,2\,) Потік може мати доступ до автоматичних змінних іншого потоку, але це небезпечно.%good

(\,3\,) Потік може мати доступ до автоматичних змінних іншого потоку і це безпечно.%bad

(\,4\,) Динамічні змінні не можуть розділятися між потоками.%bad



\begin{center}
\textbf{Завдання 5}
\end{center}
У цьому фрагменті коду
\begin{verbatim}
template <class T> void f(T a1, T a2){}

void g(){
    int i=5;
    f(unique_ptr<int>(new int(++i)), unique_ptr<int>(new int(++i)));
\end{verbatim}

(\,1\,) є невизначена поведінка%good

(\,2\,) є гонка даних%bad

(\,3\,) є стан гонки%bad

(\,4\,) можливий виток ресурсів%bad


}
\end {large}


\end{document}

